\section{Iterative solution of the flow equation}\label{sec:iterative}

In this section, we introduce the gradient flow
and derive the Feynman rules for the associated flow-line diagrams.
The pure $\SUn$ gauge theory in $D=4-2\eps$
euclidean dimensions is considered and
the fundamental gauge field $A_{\mu}(x)$ is normalized so that
its action at bare coupling $g_0$ is given by
\begin{align}
  S=-\frac{1}{2g_0^2}\int\rmd^Dx\,\tr\{F_{\mu\nu}(x)F_{\mu\nu}(x)\},
  \label{eq:2.1}\\[2pt]
  F_{\mu\nu}=\partial_{\mu}A_{\nu}-\partial_{\nu}A_{\mu}+[A_{\mu},A_{\nu}]
  \label{eq:2.2}
\end{align}
(see appendix~\ref{app:A} for unexplained notation).

\subsection{Definition of the gradient flow}\label{ssec:defflow}

The gradient flow evolves the gauge field as a function
of a parameter $t\geq0$ that is referred to as the flow time.
Starting from the fundamental gauge field,
\begin{equation}
  \left.B_{\mu}\right|_{t=0}=A_{\mu},
  \label{eq:2.3}
\end{equation}
the time-dependent field $B_{\mu}(t,x)$
is determined by the differential equation
\begin{align}
  \partial_tB_{\mu}=D_{\nu}G_{\nu\mu}+\alpha_0D_{\mu}\partial_{\nu}B_{\nu},
  \label{eq:2.4}\\[2pt]
  G_{\mu\nu}=\partial_{\mu}B_{\nu}-\partial_{\nu}B_{\mu}+[B_{\mu},B_{\nu}],
  \qquad
  D_{\mu}=\partial_{\mu}+[B_{\mu},\,\cdot\;].
  \label{eq:2.5}
\end{align}
The name ``gradient flow'' derives from the fact that first term on
the right of eq.~\eqref{eq:2.4} is proportional to the gradient of the gauge
action along the flow.
Note that neither the initial condition \eqref{eq:2.3} nor the
flow equation \eqref{eq:2.4} involve the gauge coupling.

The second term on the right of eq.~\eqref{eq:2.4} is included in order to
damp the evolution of the gauge degrees of freedom of the
field. As in the case of the Langevin equation \cite{ZinnZwanziger},
some technicalities in the perturbative analysis of the flow
can be avoided in this way without affecting the evolution
of the gauge-invariant observables.
The latter are in fact independent
of the parameter $\alpha_0$, since
the solutions of eq.~\eqref{eq:2.4} obtained at different values
of $\alpha_0$ are related by
a (time-dependent) gauge transformation \cite{WilsonFlow}.

\subsection{Expansion in powers of the fundamental gauge field}

Equation \eqref{eq:2.4} may be split into a linear and remainder part
according to
\begin{align}
  \partial_tB_{\mu}=\partial_{\nu}\partial_{\nu}B_{\mu}
  +(\alpha_0-1)\partial_{\mu}\partial_{\nu}B_{\nu}+R_{\mu},
  \label{eq:2.6}\\[2pt]
  R_{\mu}=2[B_{\nu},\partial_{\nu}B_{\mu}]-
  [B_{\nu},\partial_{\mu}B_{\nu}]+(\alpha_0-1)[B_{\mu},\partial_{\nu}B_{\nu}]
  +[B_{\nu},[B_{\nu},B_{\mu}]].
  \label{eq:2.7}
\end{align}
The linearized equation can be solved using the heat kernel
\begin{equation}
  K_t(z)_{\mu\nu}=
  \int_p{\rme^{ipz}\over p^2}\bigl\{
  (\delta_{\mu\nu}p^2-p_{\mu}p_{\nu})\rme^{-tp^2}+
  p_{\mu}p_{\nu}\rme^{-\alpha_0tp^2}\bigr\},
  \label{eq:2.8}
\end{equation}
where
\begin{equation}
  \int_p=\int{\rmd^Dp\over(2\pi)^D}.
  \label{eq:2.9}
\end{equation}
Taking the boundary condition \eqref{eq:2.3} into account, the flow equation
may then be cast in the integral form
\begin{equation}
  B_{\mu}(t,x)=\int\rmd^Dy\,\Bigl\{
  K_t(x-y)_{\mu\nu}A_{\nu}(y)
  +\int_0^t\rmd s\,K_{t-s}(x-y)_{\mu\nu}R_{\nu}(s,y)\Bigr\}.
  \label{eq:2.10}
\end{equation}
From this representation the retarded character of the equation
is evident and it is also quite
clear that the sensitivity to the initial value of the field dies away as
$t$ increases, although only slowly so at small momenta
($\alpha_0$ is assumed to be positive).

When passing to momentum space,
\begin{equation}
  B_{\mu}(t,x)=\int_p\rme^{ipx}\tilde{B}_{\mu}(t,p),
  \label{eq:2.11}
\end{equation}
the integral equation \eqref{eq:2.10} becomes
\begin{equation}
  \tilde{B}_{\mu}(t,p)=\tilde{K}_t(p)_{\mu\nu}\tilde{A}_{\nu}(p)+
  \int_0^t\rmd s\,\tilde{K}_{t-s}(p)_{\mu\nu}\tilde{R}_{\nu}(s,p).
  \label{eq:2.12}
\end{equation}
It is helpful at this point to introduce
the vertices $\vbulk{2,0}$ and $\vbulk{3,0}$ through
\begin{align}
  \tilde{R}_{\mu}^a(t,p)=\sum_{n=2}^3
  {1\over n!}\int_{q_1}\ldots\int_{q_n}(2\pi)^D\delta(p+q_1+\ldots+q_n)
  \notag\\[2.5ex]
  \times\vbulk{n,0}(p,q_1,\ldots,q_n)^{ab_1\ldots b_n}_{\mu\nu_1\ldots\nu_n}
  \tilde{B}^{b_1}_{\nu_1}(t,-q_1)\ldots\tilde{B}^{b_n}_{\nu_n}(t,-q_n)
  \label{eq:2.13}
\end{align}
and the requirement that they are totally symmetric in the
momentum-index combinations $(q_1,\nu_1,a_1),\ldots,(q_n,\nu_n,a_n)$.
The solution of the integral equation in powers of the
fundamental gauge field,
\begin{align}
  \tilde{B}_{\mu}^a(t,p)=\tilde{K}_t(p)_{\mu\nu}\tilde{A}_{\nu}^a(p)
  +\frac{1}{2}\int_0^t\rmd s\,\tilde{K}_{t-s}(p)_{\mu\nu}
  \int_{q,r}(2\pi)^D\delta(p-q-r)
  \notag\\[2.5ex]
  \times
  \vbulk{2,0}(p,-q,-r)^{abc}_{\nu\rho\sigma}
  \tilde{K}_{s}(q)_{\rho\delta}\tilde{K}_{s}(r)_{\sigma\tau}
  \tilde{A}^b_{\delta}(q)\tilde{A}^c_{\tau}(r)+\ldots,
  \label{eq:2.14}
\end{align}
is then obtained through iteration, i.e.~by
recursively inserting the equation on the right of itself.

The vertices $\vbulk{2,0}$ and $\vbulk{3,0}$ are given explicitly
in appendix~\ref{app:B}. Note that the momentum-index combination
$(p,\mu,a)$ plays a special r\^ole in eq.~\eqref{eq:2.13}. In particular,
the vertices are symmetric only in their other arguments.

\begin{figure}[t]
\centerline{\includegraphics[height=2.7cm]{images/diagram1}}
\caption{The diagrams contributing to $\tilde{B}^a_{\mu}(t,p)$
are directed tree graphs with
a single external line (a little square is drawn at the end of this line).
Flow lines always start from a one-point vertex
(circle with a cross)
or a flow vertex (circle) and end at another flow vertex
if the line is not the external one. Each flow vertex has
one outgoing flow line that corresponds to the momentum-index
combination $(p,\mu,a)$ in eq.~\eqref{eq:2.13}.}
\label{fig:trees}
\end{figure}

\subsection{Flow-line diagrams}

The terms contributing to the expansion \eqref{eq:2.14} in powers of the
fundamental gauge field can be graphically represented by
Feynman diagrams (see fig.~\ref{fig:trees}).
There are two kinds of vertices in these diagrams,
the flow vertices $\vbulk{2,0}$ and $\vbulk{3,0}$ introduced
in the previous subsection and the one-point vertex
\begin{equation}
  \raisebox{-0.45cm}{\includegraphics[width=1.0cm]{images/diagram3}}%
  \;=\;\tilde{A}^a_{\mu}(p).
  \label{eq:2.15}
\end{equation}
Each flow vertex is inserted at some flow time that
is integrated from zero to infinity, while the one-point
vertices reside at time zero.
The vertices are connected through the flow lines
\begin{equation}
  \raisebox{-0.55cm}{\includegraphics[width=2.5cm]{images/prop1}}%
  =\;\delta^{ab}\theta(t-s)\tilde{K}_{t-s}(p)_{\mu\nu},
  \label{eq:2.16}
\end{equation}
where $t$ and $s$ are the flow times at the endpoints of the line.
In view of the retarded nature
of the propagator \eqref{eq:2.16}, the flow time increases
from zero at the one-point vertices to time $t$
at the outer end of the external line as one follows
the arrows in the diagram.
In particular,
the times associated with the vertices
are effectively integrated only up to $t$ (rather than infinity).

As far the momenta, the index contractions and the symmetry factors
are concerned, the Feynman rules are the usual ones.
It is then not difficult to show that the sum of all
diagrams solves the flow equation \eqref{eq:2.12} order
by order in the fundamental gauge field.
